% FAC-EN-U0071
% Source: fac_body.tex lines 2810--2885.
% Authority: PDF physical pp. 72--73; printed pp. 267--268
% (Chapter III section 5 / no. 73 complete; hard stop before no. 74).

\medskip
\phantomsection\label{fac:III:5}
\phantomsection\label{fac:III:5:no:73}
\textbf{73. Relations between the functors
$\operatorname{Ext}_S^q$ and $\operatorname{Ext}_{\mathcal O_x}^q$.}
Let $M$ and $N$ be two graded $S$-modules. If $x$ is a point of
$X=\mathbf P_r(K)$, we defined the $\mathcal O_x$-modules $M_x$ and
$N_x$ in \hyperref[fac:III:2:no:57]{no.~57}; we shall relate the
$\operatorname{Ext}_{\mathcal O_x}^q(M_x,N_x)$ to the graded
$S$-module $\operatorname{Ext}_S^q(M,N)$.

\phantomsection\label{fac:III:5:no:73:prop:1}
\textsc{Proposition 1.} \textit{Suppose that $M$ is finitely generated.
Then:}

\phantomsection\label{fac:III:5:no:73:prop:1:a}
\textit{(a) The sheaf
$\mathfrak A(\operatorname{Hom}_S(M,N))$ is isomorphic to the sheaf
$\operatorname{Hom}_{\mathcal O}(\mathfrak A(M),\mathfrak A(N))$.}

\phantomsection\label{fac:III:5:no:73:prop:1:b}
\textit{(b) For every $x\in X$, the $\mathcal O_x$-module
$\operatorname{Ext}_S^q(M,N)_x$ is isomorphic to the
$\mathcal O_x$-module
$\operatorname{Ext}_{\mathcal O_x}^q(M_x,N_x)$.}

First define a homomorphism
\phantomsection\label{fac:III:5:no:73:map:iota-x}
\[
\iota_x:\operatorname{Hom}_S(M,N)_x
\longrightarrow\operatorname{Hom}_{\mathcal O_x}(M_x,N_x).
\]
An element of the first module is a fraction $\varphi/P$, where
$\varphi\in\operatorname{Hom}_S(M,N)_n$, $P\in S(x)$, and $P$ is
homogeneous of degree $n$. If $m/P'$ is an element of $M_x$, then
$\varphi(m)/(PP')$ is an element of $N_x$ depending only on
$\varphi/P$ and $m/P'$, and the assignment
\[
\frac{m}{P'}\longmapsto\frac{\varphi(m)}{PP'}
\]
is a homomorphism
$\iota_x(\varphi/P):M_x\longrightarrow N_x$; this defines $\iota_x$.
By \hyperref[fac:I:2:no:14:prop:5]{Proposition~5 of no.~14},
$\operatorname{Hom}_{\mathcal O_x}(M_x,N_x)$ may be identified with
\phantomsection\label{fac:III:5:no:73:formula:Hom-stalk}
\[
\operatorname{Hom}_{\mathcal O}(\mathfrak A(M),\mathfrak A(N))_x.
\]
Under this identification, $\iota_x$ becomes
\[
\iota_x:\mathfrak A(\operatorname{Hom}_S(M,N))_x
\longrightarrow
\operatorname{Hom}_{\mathcal O}(\mathfrak A(M),\mathfrak A(N))_x,
\]
and one readily checks that the collection of the $\iota_x$ is a
homomorphism
\phantomsection\label{fac:III:5:no:73:map:iota}
\[
\iota:\mathfrak A(\operatorname{Hom}_S(M,N))
\longrightarrow
\operatorname{Hom}_{\mathcal O}(\mathfrak A(M),\mathfrak A(N)).
\]

\phantomsection\label{fac:III:5:no:73:proof:free-case}
When $M$ is a finitely generated free module, $\iota_x$ is bijective:
indeed, it suffices to verify this when $M=S(n)$, in which case it is
immediate.

\phantomsection\label{fac:III:5:no:73:proof:resolution-case}
Now let $M$ be an arbitrary finitely generated graded $S$-module, and
choose a resolution of $M$:
\phantomsection\label{fac:III:5:no:73:formula:resolution}
\[
\cdots\longrightarrow L^{q+1}\longrightarrow L^q
\longrightarrow\cdots\longrightarrow L^0
\longrightarrow M\longrightarrow0,
\]
\phantomsection\label{fac:III:5:no:73:def:complex-C}
in which the $L^q$ are finitely generated free modules. Consider the
complex $C$ formed by the $\operatorname{Hom}_S(L^q,N)$. Its
cohomology groups are the $\operatorname{Ext}_S^q(M,N)$.
\phantomsection\label{fac:III:5:no:73:def:boundaries-cocycles}
In other words, if $B^q$ and $Z^q$
denote the submodules of $C^q$ consisting respectively of coboundaries
and cocycles, there are exact sequences
\phantomsection\label{fac:III:5:no:73:formula:cycles-boundaries}
\[
0\longrightarrow Z^q\longrightarrow C^q
\longrightarrow B^{q+1}\longrightarrow0
\]
and
\[
0\longrightarrow B^q\longrightarrow Z^q
\longrightarrow\operatorname{Ext}_S^q(M,N)\longrightarrow0.
\]

Since the functor $M\mapsto\mathfrak A(M)$ is exact, the sequences
\phantomsection\label{fac:III:5:no:73:formula:localized-cycles-boundaries}
\[
0\longrightarrow Z_x^q\longrightarrow C_x^q
\longrightarrow B_x^{q+1}\longrightarrow0
\]
and
\[
0\longrightarrow B_x^q\longrightarrow Z_x^q
\longrightarrow\operatorname{Ext}_S^q(M,N)_x\longrightarrow0
\]
are also exact.

But, \hyperref[fac:III:5:no:73:proof:free-case]{by what precedes}, $C_x^q$
is isomorphic to
$\operatorname{Hom}_{\mathcal O_x}(L_x^q,N_x)$. The modules
$\operatorname{Ext}_S^q(M,N)_x$ are therefore isomorphic to the
cohomology groups of the complex formed by the
$\operatorname{Hom}_{\mathcal O_x}(L_x^q,N_x)$; and since the $L_x^q$
are plainly free $\mathcal O_x$-modules, this is precisely the
definition of $\operatorname{Ext}_{\mathcal O_x}^q(M_x,N_x)$. This proves
\hyperref[fac:III:5:no:73:prop:1:b]{(b)}. For $q=0$,
\hyperref[fac:III:5:no:73:proof:resolution-case]{the preceding argument}
shows that $\iota_x$ is bijective; hence $\iota$ is an isomorphism, which
proves \hyperref[fac:III:5:no:73:prop:1:a]{(a)}.
